\documentclass{article}
\usepackage{graphicx}
\usepackage{amsthm}
\usepackage{amssymb}
\usepackage{amsmath}
\usepackage{parskip}
\usepackage[margin=.85in]{geometry}
\usepackage{tikz}
\usepackage{titling}

 % move title upward

\title{15.C57 Optimization Methods HW2}
\author{Chris Low}
\date{} % removes the date

\begin{document}

\maketitle

\section*{Problem 1}

\subsection*{(a)}

\textbf{True.}

We prove the statement by contrapositive. Let $v^*$ be the optimal primal
cost. Suppose
\[
v^* \geq c^\top x^*.
\]
Since $x^*$ is feasible and we are minimizing,
\[
v^* \leq c^\top x^*.
\]
Therefore,
\[
v^*=c^\top x^*,
\]
so $x^*$ is optimal.

Since $x^*$ may be degenerate, there may be multiple bases corresponding
to the same $x^*$, and some of these bases may not have nonnegative
reduced costs. However, because $x^*$ is already optimal, we can pivot
through these different bases until we reach one with nonnegative reduced
costs. If we pivot from a basis with a negative reduced cost, the pivot
cannot actually improve the objective since $x^*$ is already optimal.
Therefore, these pivots must be degenerate. That is, the basis changes, but the
solution remains $x^*$. Thus, we can reach a basis $B$ corresponding to
$x^*$ such that
\[
\bar c_j
=
c_j-c_B^\top B^{-1}A_j
\geq 0
\qquad \forall j.
\]
For this basis, the associated dual basic solution is
\[
p^\top=c_B^\top B^{-1}.
\]
Therefore,
\[
c_j-p^\top A_j\geq0
\qquad \forall j,
\]
which gives
\[
p^\top A_j\leq c_j
\qquad \forall j.
\]
Hence $p$ is dual feasible.

Thus, if $x^*$ is optimal, at least one basis corresponding to $x^*$
has a feasible associated dual basic solution. By contrapositive, if
every basis corresponding to $x^*$ has an infeasible associated dual
basic solution, then $x^*$ is not optimal. Since $x^*$ is feasible and
we are minimizing,
\[
v^*<c^\top x^*.
\]

Therefore, the statement is true. \qed


\newpage


\subsection*{(b)}


\textbf{True.}

The original dual problem is
\[
\max \ p^\top b
\]
subject to
\[
p^\top A \leq c^\top.
\]

If we remove the $i$th equality constraint from the primal, then the
corresponding dual variable $p_i$ is also removed. The resulting dual can
be written as
\[
\max \sum_{j\neq i} p_j b_j
\]
subject to
\[
\sum_{j\neq i} p_j A_j \leq c^\top,
\]
where $A_j$ denotes the corresponding row of $A$.

Now consider the original dual and instead set
\[
p_i=0.
\]
For the objective,
\[
p^\top b
=
\sum_{j=1}^m p_jb_j
=
\sum_{j\neq i}p_jb_j+p_ib_i
=
\sum_{j\neq i}p_jb_j.
\]
Similarly, for the dual constraint,
\[
p^\top A
=
\sum_{j=1}^m p_jA_j
=
\sum_{j\neq i}p_jA_j+p_iA_i
=
\sum_{j\neq i}p_jA_j.
\]
Therefore, setting $p_i=0$ gives us the same objective and constraints
\[
\max \sum_{j\neq i}p_jb_j
\]
subject to
\[
\sum_{j\neq i}p_jA_j\leq c^\top,
\]

as we get from removing the $i$th primal
equality constraint. Therefore, the statement is true. \qed


\vspace{5mm}
\subsection*{(c)}


\textbf{True.}

Suppose the unboundedness criterion in the primal simplex algorithm is
satisfied. Then there is some nonbasic variable $x_j$ such that
\[
\bar c_j<0
\]
and
\[
B^{-1}A_j\leq 0.
\]

If $x_j$ enters the basis, the corresponding direction for the basic
variables is
\[
d_B=-B^{-1}A_j.
\]
 
which implies
\[
d_B\geq0.
\]

Starting from the current basic feasible solution, we can move in this
direction by an amount $\theta\geq0$:
\[
x_B(\theta)
=
x_B+\theta d_B
=
x_B-\theta B^{-1}A_j.
\]
Since $x_B\geq0$ and $d_B\geq0$, we have
\[
x_B(\theta)\geq0
\]
for every $\theta\geq0$. Therefore, $\theta$ can be increased
arbitrarily while the primal solution remains feasible.

The objective value along this direction is
\[
c^\top x(\theta)
=
c^\top x+\theta\bar c_j.
\]
Since $\bar c_j<0$, as $\theta\to\infty$,
\[
c^\top x(\theta)\to-\infty.
\]
Thus, the primal problem is unbounded below.

Now suppose, for contradiction, that the dual is feasible, and we let $p$
be a dual feasible solution. Then
\[
p^\top A\leq c^\top.
\]
Since $x(\theta)\geq0$, multiplying by $x(\theta)$ gives us 
\[
p^\top A x(\theta)\leq c^\top x(\theta).
\]
And because we know now that $x(\theta)$ is primal feasible,
\[
Ax(\theta)=b,
\]
so
\[
p^\top b\leq c^\top x(\theta).
\]

However, $p^\top b$ is fixed while $
c^\top x(\theta)\to-\infty.
$
For sufficiently large $\theta$, we would therefore have
\[
c^\top x(\theta)<p^\top b,
\]
contradicting the inequality above. Hence no dual feasible solution can
exist, so the dual problem is infeasible. \qed



\vspace{5mm}
\subsection*{(d)}

\textbf{True.}

If $x^*$ is a BFS, then every variable with positive flow must be a
basic variable, since all nonbasic variables are equal to zero. The columns corresponding to the basic variables form the basis matrix \(B\) so they must be linearly independent for \(B\) to be invertible.


Suppose the positive-flow arcs contain an undirected cycle. The columns
corresponding to the arcs in a cycle are linearly dependent. For example, for the cycle $1\to2\to3\to1$,
\[
A_{12}+A_{23}+A_{31}
=
\begin{bmatrix}
-1\\
1\\
0
\end{bmatrix}
+
\begin{bmatrix}
0\\
-1\\
1
\end{bmatrix}
+
\begin{bmatrix}
1\\
0\\
-1
\end{bmatrix}
=
\begin{bmatrix}
0\\
0\\
0
\end{bmatrix}.
\]
More generally, if some arcs point in the opposite direction around the
cycle, we can add or subtract their columns so that the entries at each
node cancel, again giving a nonzero linear combination equal to zero.

Therefore, the positive-flow arcs in the cycle are linearly dependent and cannot all be part of a basis. However, since they all have positive flow, they would all have to be basic if \(x^*\) were a BFS. This would require the same columns to be both linearly independent (because they are basic) and linearly dependent (because they form a cycle), which is a contradiction. Therefore, \(x^*\) cannot be a basic feasible solution. \qed


\newpage




\subsection*{(e)}

\textbf{False.}

The basic arcs of a BFS form a spanning tree, but the arcs with strictly
positive flow do not necessarily form a spanning tree. This is because
a BFS can be degenerate, meaning that some basic variables can be equal
to zero. Therefore, some arcs can be part of the basis but carry zero
flow.

For example, consider a network with nodes $1,2,3$ and arcs
\[
(1,2), \qquad (2,3),
\]
with supplies
\[
b_1=1,\qquad b_2=-1,\qquad b_3=0.
\]
The two arcs form a spanning tree and can therefore form a basis. The
corresponding feasible flow is
\[
x_{12}=1,\qquad x_{23}=0.
\]
Thus, this is a degenerate BFS since the basic variable $x_{23}$ is
equal to zero.

However, the only arc carrying strictly positive flow is $(1,2)$. This
arc does not connect node $3$, so the positive-flow arcs do not form a
spanning tree. Therefore, the statement is false.


\newpage


\section*{Problem 2}

\subsection*{(a)}


Introduce a variable $z$ to represent the maximum term. The original
problem can be written as the LP
\[
\begin{aligned}
\min_{x,z}\quad & z\\
\text{s.t.}\quad & a_i^\top x-b_i\leq z,
\qquad i=1,\ldots,m.
\end{aligned}
\]

As given, we let $p\geq0$ satisfy
\[
p^\top A=0^\top
\qquad\text{and}\qquad
\sum_{i=1}^m p_i=1.
\]

We then multiply the $i$th constraint by $p_i\geq0$ and sum them over all $i$ which gives us
\[
\sum_{i=1}^m p_i(a_i^\top x-b_i)
\leq
z\sum_{i=1}^m p_i.
\]

We can see the left-hand side becomes
\[
p^\top Ax-p^\top b
=
-p^\top b,
\]
because $p^\top A=0^\top$, while the right-hand side becomes
\[
z\sum_{i=1}^m p_i=z.
\]

Therefore,
\[
-p^\top b\leq z
\]
for every feasible $(x,z)$. Since the minimum value of $z$ is $v$, we
have
\[
-p^\top b\leq v.
\]




\vspace{5mm}
\subsection*{(b)}


From part (a), every feasible $p$ satisfies
$
-p^\top b \leq v.
$
We now show that the best such lower bound is exactly $v$.

From part (a), the original problem can be written as
\[
\begin{aligned}
\min_{x,z}\quad & z\\
\text{s.t.}\quad & a_i^\top x-b_i\leq z,
\qquad i=1,\ldots,m.
\end{aligned}
\]

We multiply the constraints by $-1$\[
-a_i^\top x+z\geq -b_i,
\]
so in matrix form the primal is
\[
\begin{aligned}
\min_{x,z}\quad & z\\
\text{s.t.}\quad & -Ax+\mathbf{1}z\geq -b,
\end{aligned}
\]
where $x$ and $z$ are unrestricted.

Let $p_i\geq 0$ be the dual variable to the $i$th
primal constraint. Since each $x_j$ is unrestricted, its corresponding
dual constraint is an equality. And because the objective coefficient of
each $x_j$ is zero (we only have $z$ in the objective above), these constraints give
\[
-A^\top p=0,
\]
or equivalently,
\[
p^\top A=0^\top.
\]

The variable $z$ is also unrestricted. Its coefficient in every primal
constraint is $1$, and its objective coefficient is $1$. Therefore, its
corresponding dual constraint is
\[
\mathbf{1}^\top p=1,
\]
or
\[
\sum_{i=1}^m p_i=1.
\]

Finally, the right-hand side of the primal constraints is $-b$, so the
dual objective is
\[
\max_p -p^\top b.
\]

Therefore, the dual problem is
\[
\begin{aligned}
\max_p\quad & -p^\top b\\
\text{s.t.}\quad & p^\top A=0^\top,\\
& \sum_{i=1}^m p_i=1,\\
& p\geq 0.
\end{aligned}
\]

This is exactly the optimization problem given in part (b). Since the
primal has finite optimal value $v$, strong duality implies that the
dual has the same optimal value. Hence,
\[
\max_p\{-p^\top b:
p^\top A=0^\top,\ \mathbf{1}^\top p=1,\ p\geq0\}=v.
\]
\qed


\newpage


\section*{Problem 3}


\subsection*{(a)}


The objective function
\[
\sum_{t=1}^T \alpha^{t-1}(d_t x_t-cy_t)
+\alpha^T d_{T+1}z_T
\]
maximizes the total discounted profit over the planning horizon. In
period $t$, $d_t x_t$ is the revenue from selling $x_t$ units at price
$d_t$, while $cy_t$ is the transportation cost of sending $y_t$ units
to long-term storage. These are discounted by $\alpha^{t-1}$ since
future revenues and costs are valued less than current ones. Because retrieval
from storage is free, we do not include $w_t$ in the objective. Finally,
any inventory $z_T$ remaining at the end of the horizon is valued at the
salvage price $d_{T+1}$, and we take into the discount rate as well. This gives us the discounted terminal value
$\alpha^T d_{T+1}z_T$.

The first constraint
\[
x_t+y_t-w_t=b_t
\]
can be rewritten as
\[
x_t+y_t=b_t+w_t.
\]
This is the product flow balance in period $t$. The available product
consists of new production $b_t$ and units retrieved from storage $w_t$.
All available product must either be sold, $x_t$, or sent to long-term
storage, $y_t$.

The second constraint
\[
z_t+w_t-z_{t-1}-fy_t=0
\]
can be rewritten as
\[
z_t=z_{t-1}+fy_t-w_t.
\]
This is the inventory balance for long-term storage. The end-of-period
inventory at $t$ equals the inventory carried over from the previous period,
plus the fraction $f$ of the newly stored product that survives, minus
the amount $w_t$ retrieved from storage.

The condition $
z_0=0
$
means that long-term storage is empty at the beginning of the planning
horizon.

Finally,
$
x_t,y_t,w_t,z_t\geq0
$
because the quantities sold, sent to storage, retrieved, and kept in
storage cannot be negative. Since demand is unlimited, we do not have an
upper bound imposed on $x_t$.


\vspace{5mm}
\subsection*{(b)}

Let $p_t$ and $q_t$ be the dual variables associated with the first and
second equality constraints, respectively:
\[
x_t+y_t-w_t=b_t,
\]
and
\[
z_t+w_t-z_{t-1}-fy_t=0.
\]
Because both primal constraints are equalities, the dual variables
$p_t$ and $q_t$ are unrestricted in sign.

The primal is a maximization problem and
$
x_t,y_t,w_t,z_t\geq 0,
$
so each primal variable gives a ``greater than or equal to'' constraint
in the dual.

For $x_t$, its coefficient is $1$ in the first constraint and $0$ in
the second constraint. Its objective coefficient is
$\alpha^{t-1}d_t$. Therefore,
\[
p_t\geq \alpha^{t-1}d_t,
\qquad t=1,\ldots,T.
\]

For $y_t$, its coefficients in the first and second constraints are
$1$ and $-f$, respectively. Its objective coefficient is
$-\alpha^{t-1}c$. Therefore,
\[
p_t-fq_t\geq -\alpha^{t-1}c,
\qquad t=1,\ldots,T.
\]

For $w_t$, its coefficients in the first and second constraints are
$-1$ and $1$, respectively. Because $w_t$ does not appear in the
objective, its objective coefficient is zero. Therefore,
\[
-p_t+q_t\geq 0,
\qquad t=1,\ldots,T.
\]

For $z_t$, where $t=1,\ldots,T-1$, it appears with coefficient $1$ in
the period-$t$ storage constraint and coefficient $-1$ in the
period-$(t+1)$ storage constraint. And because $z_t$ does not appear in the
objective for $t<T$, we obtain
\[
q_t-q_{t+1}\geq 0,
\qquad t=1,\ldots,T-1.
\]

Finally, we need to take into account of the last period $z_T$, which appears only in the period-$T$ storage constraint with
coefficient $1$, and its objective coefficient is
$\alpha^T d_{T+1}$. Hence,
\[
q_T\geq \alpha^T d_{T+1}.
\]

The dual objective is obtained from the right-hand sides of the primal
constraints. The first set of constraints has right-hand side $b_t$,
while the second set has right-hand side zero. Thus, the dual objective
is
\[
\min \sum_{t=1}^T b_t p_t.
\]

Therefore, the dual problem is
\[
\begin{aligned}
\min_{p,q}\quad
& \sum_{t=1}^T b_t p_t\\
\text{s.t.}\quad
& p_t\geq \alpha^{t-1}d_t,
&& t=1,\ldots,T,\\
& p_t-fq_t\geq -\alpha^{t-1}c,
&& t=1,\ldots,T,\\
& -p_t+q_t\geq 0,
&& t=1,\ldots,T,\\
& q_t-q_{t+1}\geq 0,
&& t=1,\ldots,T-1,\\
& q_T\geq \alpha^T d_{T+1},\\
& p_t,q_t\text{ unrestricted in sign},
&& t=1,\ldots,T.
\end{aligned}
\]

\vspace{5mm}
\subsection*{(c)}


We first show that the proposed solution is dual feasible. Recall that
the dual constraints can be written as
\[
p_t \geq \alpha^{t-1}d_t,
\]
\[
p_t \geq fq_t-\alpha^{t-1}c,
\]
\[
q_t \geq p_t,
\]
\[
q_t \geq q_{t+1},
\qquad t=1,\ldots,T-1,
\]
\[
q_T \geq \alpha^T d_{T+1}.
\]

By the proposed definition,
\[
p_t
=
\max\left\{
fq_t-\alpha^{t-1}c,\,
\alpha^{t-1}d_t
\right\}.
\]
Since a maximum is at least each of its terms, we immediately have
\[
p_t\geq fq_t-\alpha^{t-1}c
\]
and
\[
p_t\geq\alpha^{t-1}d_t.
\]
Thus, the first two dual constraints are satisfied.

For $t=1,\ldots,T-1$, the proposed definition is
\[
q_t
=
\max\left\{
q_{t+1},\,
\alpha^{t-1}d_t
\right\}.
\]
Therefore,
\[
q_t\geq q_{t+1}
\]
and
\[
q_t\geq\alpha^{t-1}d_t.
\]

At the final period,
\[
q_T
=
\max\left\{
\alpha^T d_{T+1},\,
\alpha^{T-1}d_T
\right\},
\]
so
\[
q_T\geq\alpha^T d_{T+1}
\]
and
\[
q_T\geq\alpha^{T-1}d_T.
\]
Thus, the constraints involving $q_t\geq q_{t+1}$ and
$q_T\geq\alpha^T d_{T+1}$ are satisfied.

It remains only to verify
\[
q_t\geq p_t.
\]
From the definitions of $q_t$ and $q_T$, we have for every
$t=1,\ldots,T$,
\[
q_t\geq\alpha^{t-1}d_t.
\]
Also, since $0<f<1$, $c\geq0$, and $q_t\geq0$, we have 
\[
fq_t-\alpha^{t-1}c
\leq fq_t
\leq q_t.
\]
Therefore, $q_t$ is at least both terms in
\[
p_t
=
\max\left\{
fq_t-\alpha^{t-1}c,\,
\alpha^{t-1}d_t
\right\}.
\]
So we conclude that 
\[
q_t\geq p_t.
\]

Thus, all of the dual constraints are satisfied, so the proposed
solution is \textbf{dual feasible}.



\vspace{3mm}
We now show that these choices are \textbf{optimal} by working backward from
period $T$. Any feasible solution must satisfy (as given)
\[
q_T\geq\alpha^T d_{T+1}.
\]
Also, since we have
\[
q_T\geq p_T
\qquad\text{and}\qquad
p_T\geq\alpha^{T-1}d_T,
\]
we must also have
\[
q_T\geq\alpha^{T-1}d_T.
\]
Therefore every feasible solution must satisfy
\[
q_T\geq
\max\left\{
\alpha^T d_{T+1},\,
\alpha^{T-1}d_T
\right\}.
\]
Thus the proposed value
$
q_T=
\max\left\{
\alpha^T d_{T+1},\,
\alpha^{T-1}d_T
\right\}
$
is the smallest feasible value of $q_T$.

Once $q_T$ is chosen, feasibility requires
\[
p_T\geq fq_T-\alpha^{T-1}c
\]
and
\[
p_T\geq\alpha^{T-1}d_T.
\]
Hence the smallest feasible value is
\[
p_T=
\max\left\{
fq_T-\alpha^{T-1}c,\,
\alpha^{T-1}d_T
\right\}.
\]

We can now repeat the same argument backward. For any
$t=1,\ldots,T-1$, feasibility requires
\[
q_t\geq q_{t+1}
\]
and, since
\[
q_t\geq p_t\geq\alpha^{t-1}d_t,
\]
also
\[
q_t\geq\alpha^{t-1}d_t.
\]
Therefore the smallest feasible value of $q_t$ is
\[
q_t=
\max\left\{
q_{t+1},\,
\alpha^{t-1}d_t
\right\}.
\]

Once $q_t$ is determined, the two lower bounds on $p_t$ give
\[
p_t=
\max\left\{
fq_t-\alpha^{t-1}c,\,
\alpha^{t-1}d_t
\right\}
\]
as the smallest feasible value of $p_t$.

Thus, working backward from $T$ to $1$, the proposed construction
chooses the smallest feasible values of $q_t$ and $p_t$ in every
period. Since the dual objective is
$
\min \sum_{t=1}^T b_t p_t
$
and $b_t\geq0$, choosing the smallest feasible $p_t$ minimizes the
objective. Therefore, the proposed solution is optimal. \qed






\vspace{5mm}
\subsection*{(d)}


From part (b), the dual constraints corresponding to the primal variables are
\[
p_t\geq \alpha^{t-1}d_t,
\]
\[
p_t-fq_t\geq -\alpha^{t-1}c,
\]
\[
q_t-p_t\geq 0,
\]
\[
q_t-q_{t+1}\geq 0,
\qquad t=1,\ldots,T-1,
\]
\[
q_T\geq \alpha^T d_{T+1}.
\]

By complementary slackness, if a primal variable is positive, then its
corresponding dual constraint must be tight. Thus,
\[
x_t>0
\quad\Rightarrow\quad
p_t=\alpha^{t-1}d_t,
\]
\[
y_t>0
\quad\Rightarrow\quad
p_t=fq_t-\alpha^{t-1}c,
\]
\[
w_t>0
\quad\Rightarrow\quad
q_t=p_t,
\]
and for $t<T$,
\[
z_t>0
\quad\Rightarrow\quad
q_t=q_{t+1},
\]
and lastly we have
\[
z_T>0
\quad\Rightarrow\quad
q_T=\alpha^T d_{T+1}.
\]

We first compute the optimal dual variables backward using part (c):
\[
q_T=\max\left\{\alpha^T d_{T+1},\alpha^{T-1}d_T\right\},
\]
\[
p_T=\max\left\{fq_T-\alpha^{T-1}c,\alpha^{T-1}d_T\right\},
\]
and for $t=T-1,\ldots,1$,
\[
q_t=\max\left\{q_{t+1},\alpha^{t-1}d_t\right\},
\]
\[
p_t=\max\left\{fq_t-\alpha^{t-1}c,\alpha^{t-1}d_t\right\}.
\]

We then reconstruct the primal solution forward from $t=1$ to $T$.

\textbf{First Case:}


Suppose
\[
\alpha^{t-1}d_t>q_{t+1}.
\]
Then
\[
q_t=\alpha^{t-1}d_t>q_{t+1}.
\]
Therefore the dual constraint corresponding to $z_t$ is slack, so by
complementary slackness,
\[
z_t=0.
\]
Since $q_t=\alpha^{t-1}d_t$, we also have
\[
p_t=\alpha^{t-1}d_t=q_t,
\]
so the dual constraints corresponding to $x_t$ and $w_t$ are tight.
Under nondegeneracy, this gives
\[
x_t>0,\qquad w_t>0,
\]
while the $y_t$ constraint is slack, so
\[
y_t=0.
\]
Using the primal balance equations,
\[
z_t=z_{t-1}+fy_t-w_t,
\]
we obtain
\[
w_t=z_{t-1},
\]
and from
\[
x_t+y_t-w_t=b_t,
\]
we obtain
\[
x_t=b_t+z_{t-1}.
\]
Hence in this case,
\[
\boxed{
w_t=z_{t-1},\qquad
y_t=0,\qquad
x_t=b_t+z_{t-1},\qquad
z_t=0.}
\]

It essentially means we retrieve and sell everything. 


\vspace{3mm}
\textbf{Second Case:}


Suppose instead
\[
q_{t+1}>\alpha^{t-1}d_t.
\]
Then
\[
q_t=q_{t+1},
\]
so the dual constraint corresponding to $z_t$ is tight and therefore,
under nondegeneracy,
\[
z_t>0.
\]
Also, $q_t>p_t$, so the dual constraint corresponding to $w_t$ is
slack, which implies
\[
w_t=0.
\]

There are now two subcases.

\newpage
\textbf{First Subcase:}


If
\[
fq_t-\alpha^{t-1}c>\alpha^{t-1}d_t,
\]
then
\[
p_t=fq_t-\alpha^{t-1}c.
\]
Thus the dual constraint corresponding to $y_t$ is tight, while the one
corresponding to $x_t$ is slack. Hence
\[
y_t>0,\qquad x_t=0.
\]
Since $w_t=0$, the first primal balance ($x_t + y_t - w_t = b_t$) gives
\[
y_t=b_t,
\]
and $z_t + w_t - z_{t-1} - fy_t = 0$ gives
\[
z_t=z_{t-1}+fb_t.
\]
Thus
\[
\boxed{
w_t=0,\qquad
y_t=b_t,\qquad
x_t=0,\qquad
z_t=z_{t-1}+fb_t.}
\]

It means we are storing production.

\vspace{3mm}

\textbf{Second Subcase:}



If
\[
\alpha^{t-1}d_t>fq_t-\alpha^{t-1}c,
\]
then
\[
p_t=\alpha^{t-1}d_t.
\]
Therefore the dual constraint corresponding to $x_t$ is tight and the
one corresponding to $y_t$ is slack. Hence
\[
x_t>0,\qquad y_t=0.
\]
Since $w_t=0$, the first primal balance constraint gives us
\[
x_t=b_t,
\]
and the storage balance gives
\[
z_t=z_{t-1}.
\]
Thus
\[
\boxed{w_t=0,\qquad
y_t=0,\qquad
x_t=b_t,\qquad
z_t=z_{t-1}.}
\]


\vspace{3mm}


\textbf{Summary of the algorithm.}


Therefore, the back-propagation algorithm proceeds in two steps.

First, starting at $t=T$, compute $q_T$ and $p_T$ using the terminal
formulas from part (c). Then move backward from $t=T-1$ to $t=1$,
computing $q_t$ and $p_t$ recursively. At each period, record which of
the three cases above applies.

After completing this backward pass, set
$
z_0=0.
$

Then move forward from $t=1$ to $t=T$. In each period, use the case
identified during the backward pass together with the primal balance
equations to determine $x_t,y_t,w_t,$ and $z_t$.

Thus, the backward pass determines the optimal action in each period,
while the forward pass determines the corresponding quantities. By
complementary slackness, the resulting primal solution is optimal.


\newpage




\section*{Problem 4.}

\subsection*{(a)}

Let \(x_{ij} \geq 0\) denote the fraction of job \(i\) processed on machine \(j\). The linear optimization problem is

\[
\begin{aligned}
\min \quad
& \sum_{i=1}^{1000}\sum_{j=1}^{20} c_{ij}x_{ij} \\
\text{s.t.}\quad
& \sum_{j=1}^{20} x_{ij} = 1,
&& i=1,\ldots,1000, \\
& \sum_{i=1}^{1000} x_{ij} \leq C_j,
&& j=1,\ldots,20, \\
& x_{ij} \geq 0,
&& i=1,\ldots,1000,\quad j=1,\ldots,20.
\end{aligned}
\]

The first constraint ensures that each job is fully processed. The
second constraint ensures that machine \(j\) processes at most \(C_j\) units of jobs.

This problem can be represented as a minimum-cost network flow problem.
We can create one node for each job, one node for each machine, and one sink
node \(s\). Each job node \(i\) has supply \(1\), each machine node has supply/demand
\(0\), and the sink has demand \(1000\). For every job \(i\) and machine
\(j\), we can create an arc from job \(i\) to machine \(j\) with cost \(c_{ij}\)
and capacity \(1\). The flow on this arc is \(x_{ij}\). For every machine \(j\), we can create an arc from machine \(j\) to the sink
with cost \(0\) and capacity \(C_j\). Thus the network has the form

\[
\boxed{\text{Job } i}
\;\xrightarrow[\text{capacity }1]{\text{cost }c_{ij}}\;
\boxed{\text{Machine } j}
\;\xrightarrow[\text{capacity }C_j]{\text{cost }0}\;
\boxed{\text{Sink}}
\]

where the job-to-machine arc has cost \(c_{ij}\) and capacity \(1\), and
the machine-to-sink arc has cost \(0\) and capacity \(C_j\).

We need flow balance at job \(i\) which gives us 
$
\sum_{j=1}^{20} x_{ij} = 1.
$
We also need flow balance at machine \(j\), which together with the capacity of its arc to
the sink gives us $
\sum_{i=1}^{1000} x_{ij} \leq C_j.
$ Therefore, this network flow formulation is equivalent to the original
linear optimization problem.


\vspace{3mm}
For a minimum-cost network flow problem, the integrality theorem states
that if all supplies, demands, and arc capacities are integers, then
there exists an optimal solution in which all flows are integers.

In our network, each job has supply $1$, the sink has demand $1000$,
and the arc capacities are $1$ and $C_j$. Assuming that each $C_j$ is
an integer, all of these quantities are integers. Therefore, by the
integrality theorem, there exists an optimal integral flow. In particular, the flows $x_{ij}$ are integers. Since
$
x_{ij} \geq 0
$
and
$\sum_{j=1}^{20} x_{ij} = 1
$
for each job $i$, we must have
$
x_{ij} \in \{0,1\}.
$ Thus, there exists an optimal solution in which each job is processed
entirely on a single machine, even though the model allows jobs to be
split across multiple machines.

\newpage

\subsection*{(c)}

\textbf{(Note: Parts (b), (e), and (f) are attached to the end of the document.)
}


From part (a), the primal problem is
\[
\begin{aligned}
\min \quad &
\sum_{i=1}^{1000}\sum_{j=1}^{20} c_{ij}x_{ij}\\
\text{s.t.}\quad
&\sum_{j=1}^{20}x_{ij}=1,
&& i=1,\ldots,1000,\\
&\sum_{i=1}^{1000}x_{ij}\leq C_j,
&& j=1,\ldots,20,\\
&x_{ij}\geq0,
&& \forall i,j.
\end{aligned}
\]

We let $p_i$ be the dual variable associated with the job constraint
$
\sum_{j=1}^{20}x_{ij}=1,
$

and let $q_j$ be the dual variable associated with the machine capacity
constraint
$
\sum_{i=1}^{1000}x_{ij}\leq C_j.
$

Since the job constraints are equalities, $p_i$ is free in sign. Since
the machine capacity constraints are $\leq$ constraints in a minimization
problem,
$
q_j\leq0.
$

The right-hand sides of the primal constraints are $1$ and $C_j$, so
the dual objective is
\[
\max \left\{
\sum_{i=1}^{1000}p_i+
\sum_{j=1}^{20}C_jq_j
\right\}.
\]

For the primal variable, each of $x_{ij}$ appears with coefficient $1$ in the job
$i$ constraint and coefficient $1$ in the machine $j$ constraint.
Since $x_{ij}\geq0$, the corresponding dual constraint is
$
p_i+q_j\leq c_{ij}.
$

Therefore, the dual problem is
\[
\boxed{
\begin{aligned}
\max \quad &
\sum_{i=1}^{1000}p_i+
\sum_{j=1}^{20}C_jq_j\\
\text{s.t.}\quad
&p_i+q_j\leq c_{ij},
&& i=1,\ldots,1000,\quad j=1,\ldots,20,\\
&p_i\text{ free},
&& i=1,\ldots,1000,\\
&q_j\leq0,
&& j=1,\ldots,20.
\end{aligned}}
\]


\textbf{Interpretation of the dual variables.} The dual variable $p_i$ is associated with the constraint that job $i$
must be fully processed. It can be interpreted as the marginal energy
cost of increasing the amount of job $i$ that needs to be processed.

The dual variable $q_j$ is associated with the capacity of machine $j$.
Since $q_j\leq 0$, a negative value means that increasing the capacity
of machine $j$ would decrease the minimum total energy consumption. If
$q_j=0$, then increasing the capacity of that machine would not improve
the optimal objective value.

The dual constraint
$
p_i+q_j\leq c_{ij}
$
can also be written as
$
p_i\leq c_{ij}-q_j.
$
Thus, $c_{ij}-q_j$ can be viewed as the energy cost of assigning job
$i$ to machine $j$, adjusted for how valuable the capacity of machine
$j$ is. Since the dual maximizes $p_i$, $p_i$ will be limited by the
machine with the smallest adjusted cost.

\vspace{3mm}

\textbf{Complementary slackness.} For each $x_{ij}$, complementary slackness gives
\[
x_{ij}\left(c_{ij}-p_i-q_j\right)=0.
\]
Therefore, if job $i$ is assigned to machine $j$, so that $x_{ij}>0$,
then
$
p_i+q_j=c_{ij}.
$
Equivalently,
$
p_i=c_{ij}-q_j.
$


Since the dual constraints require
$
p_i\leq c_{ij}-q_j
$ for every machine $j$, and the dual maximizes $p_i$, we have
$
p_i=\min_j\{c_{ij}-q_j\}.
$
Therefore, job $i$ can only be assigned to a machine that minimizes
$c_{ij}-q_j$. If
$
p_i+q_j<c_{ij},
$
then $x_{ij}=0$, so job $i$ is not assigned to machine $j$.

For the machine capacity constraints, complementary slackness gives
\[
q_j\left(C_j-\sum_{i=1}^{1000}x_{ij}\right)=0.
\]
Thus, if
$
q_j<0,$
then
$
\sum_{i=1}^{1000}x_{ij}=C_j,
$
so machine $j$ is operating at full capacity. Conversely, if
$
\sum_{i=1}^{1000}x_{ij}<C_j,
$
then $q_j=0$. Therefore, a machine with unused capacity has zero
marginal value for additional capacity.


\vspace{5mm}


\subsection*{(d)}

Let $x_{ij}\geq 0$ be the fraction of job $i$ processed on machine $j$.
Since processing all of job $i$ on machine $j$ uses $r_{ij}$ units of
utilization, processing a fraction $x_{ij}$ uses $r_{ij}x_{ij}$ units.

\[
\begin{aligned}
\min \quad
& \sum_{i=1}^{1000}\sum_{j=1}^{20} c_{ij}x_{ij} \\
\text{s.t.}\quad
& \sum_{j=1}^{20}x_{ij}=1,
&& i=1,\ldots,1000,\\
& \sum_{i=1}^{1000}r_{ij}x_{ij}\leq U_j,
&& j=1,\ldots,20,\\
& x_{ij}\geq0,
&& \forall i,j.
\end{aligned}
\]

This is not a standard network flow problem. In part (a), each job used
the same amount of machine capacity, so the amount of flow into a machine
directly represented the amount of capacity being used. Here, each job
can use a different amount of a machine's utilization because of
$r_{ij}$. Therefore, we cannot represent the utilization constraint using
a normal capacity constraint on a network arc.

The solution is also not guaranteed to be integral. For example, suppose
there is one job and two machines, where
\[
c_{11}=0,\qquad c_{12}=1,
\]
\[
r_{11}=2,\qquad r_{12}=1,
\]
and
\[
U_1=U_2=1.
\]
Then
\[
x_{11}+x_{12}=1
\]
and the utilization constraint for machine 1 is
\[
2x_{11}\leq 1.
\]
Since machine 1 has the lower cost, we would want to assign as much of
the job to machine 1 as possible. This gives
\[
x_{11}=\frac{1}{2},\qquad x_{12}=\frac{1}{2}.
\]
Thus, the optimal solution can split a job between machines. Therefore,
unlike in part (a), we are not guaranteed to have an integral optimal
solution where each job is assigned to only one machine.

\end{document}
